\section{This is the new stuff}


\begin{lem}
  Let $G$ be an $n$-vertex graph and let $\mathcal{T}=(T,(B_x\mid x\in V(T)))$ be a tree-decomposition of $G$.  Then there exists $S\subseteq V(T)$ with $|S| \le 2b-{?}$ such that each component of $G-\bigcup_{x\in S}B_x$ has less than $n/b$ vertices and each component of $T-S$ has is adjacent to at most two vertices of $S$.
\end{lem}

\begin{proof}
  Treat $T$ as a rooted tree. If $n < b$, then there is nothing to do.  Otherwise, let $x\in V(T)$ be such that $\omega(x) \ge n/b$ but $\omega(y) < n/b$ for each $y\in V(T_x)$.  Then at most one of the trees in $T-x$ has weight greater than $n/b$. Recurse on this tree, which has weight at most $n-n/b$.  This gives a set $S_0$ of size at most $b-1$ such that each component of $G-(\bigcup_{x\in S_0}B_x)$ has size less than $n/b$.  

  Let $T_0$ be the minimal subtree of $T$ that spans $S_0$ and let $S_1$ be the set of nodes in $T_0$ of degree at least three.  Then $S:=S_0\cup S_1$ has the desired properties.
\end{proof}

\begin{lem}
  Let $G$ be an $n$-vertex graph and let $\mathcal{T}=(T,(B_x\mid x\in V(T)))$ be a tree-decomposition of $G$.  Then $G$ has a tree-decomposition $\mathcal{Q}:=(Q,(C_q\mid q\in V(Q)))$ with the following properties:
  \begin{enumerate}[nosep,nolistsep,label=\rm(\alph*)]
    \item $\height(Q)\le \log_b n$;
    \item for each $q\in V(Q)$, the torso $\torso{G}{\mathcal{Q},C_q}$ has a tree-decomposition $\mathcal{T}(q):=(T(q),(D_x:x\in V(T(q))))$ such that
    \begin{enumerate}[nosep,nolistsep,label=\rm(\roman*)]
      \item $|V(T(q))|\le {?}b$;
      \item each adhesion of $\mathcal{T}(q)$ is contained in the union of at most two adhesions of $\mathcal{T}$;
      \item each bag of $\mathcal{T}(q)$ is contained in the union of at most one bag of $\mathcal{T}$ and at most two adhesions of $\mathcal{T}$.
    \end{enumerate}
  \end{enumerate}  
\end{lem}

\begin{proof}
  The proof is by induction on $n$, with a slightly more sophisticated inductive hypothesis. An inductive instance includes a graph $G$, a tree-decomposition $\mathcal{T}:=(T,(B_x\mid x\in V(T)))$, and two (not necessarily distinct) leaves $\alpha$ and $\beta$ of $T$.  In addition to providing a tree-decomposition $\mathcal{Q}$ of $G$ satisfying the conditions of the lemma, it also ensures 
  \begin{enumerate}[nosep,nolistsep,label=\rm(\roman*)]\setcounter{enumi}{3}
    \item\label{adhesion_pairs} some bag of the root decomposition $\mathcal{T}(q_0)$ contains the parent adhesions $A_\alpha$ and $A_\beta$ of $\alpha$ and $\beta$.
  \end{enumerate}
  The induction is on the number of vertices in $V(G-(B_\alpha\cup B_\beta))$.  To see that this enough to establish the original lemma, we can attach two leaves $\alpha, \beta$ to $T$, set $B_\alpha:=B_\beta:=\emptyset$, and apply the refined result.

  If $|V(G)\setminus (B_\alpha\cup B_\beta)|<b$, then then we can take $\mathcal{Q}$ to be a trivial tree-decomposition with a single bag $B_{q_0}$.  Removing redundant nodes from $T$ produces a tree-decomposition with at most $b+2\le 2b$ nodes, and adding $B_\alpha$ to each bag on the path in $T$ from $\alpha$ to $\beta$ gives a tree-decomposition of $G[B_{q_0}]=\torso{G}{\mathcal{Q},B_{q_0}}$ that satisifies the requirements of the lemma.
  % \pat{Handle base case.}

  Now assume that $n>b$.
  Let $S_0$ be a set of at most $b-1$ nodes of $T$ with the property that, for each component $C$ of $T-S$, $|\bigcup_{x\in V(C)}B_x\setminus (\bigcup_{x\in S_0} B_x)|\le n/b$.  Let $T_0$ be the minimum subtree of $T$ that spans $S_0\cup\{\alpha,\beta\}$. Let $S_1$ be the set of nodes in $S_0$ that have degree at least $3$.  Then $|S_1|\le |S_0|\le b-1$. Let $S:=S_0\cup S_1\cup\{\alpha,b\}$, so $|S|\le 2b$.

  The tree $Q$ will have root $q_0$ with $C_{q_0}:=\bigcup_{x\in S}B_x$.  We now describe the tree-decomposition $\mathcal{Q}(q_0):=(T(q_0), D_x:x\in V(T(q_0)))$.  Create the tree $Q(q_0)$ as follows.  Say that a pair $x,y\in S$ is \defin{consecutive} if the path in $T$ from $x$ to $y$ has no internal node in $S$. For each pair of consecutive nodes $x,y\in S$, introduce a vertex $v(xy)$ and add the path $x,v(xy),y$ to $Q$.  For each $x\in S_0\cup\{x,y\}$, set $D_x:=B_x$.

  Let $C$ be a component of $T-S$ and let $\{\alpha',\beta'\}=N_T(V(C))$, so $\alpha'$ and $\beta'$ are consecutive vertices of $S$. (Possibly $\alpha'=\beta'$ if $|N_T(V(C))|=1$.)  Set $D_{v(\alpha'\beta')}:=(B_x\cap B_{\alpha'})\cup(B_y\cap B_{\beta'})$.
  Note that this ensures that some bag of $\mathcal{T}(q_0)$ contains $N_G(V(C))\cap C_{q_0}$.

  Let $T':=T[V(C)\cup N_T(V(C))]$.  We apply induction on the graph $G_{\alpha'\beta'}:=G[\bigcup_{x\in V(T')}B_x]$ with the tree-decomposition $\mathcal{T}':=(T',(B_x:x\in V(T')))$ and special leaves $\alpha'$ and $\beta'$.  Since $G_{\alpha'\beta'}-(B_{\alpha'}\cup B_{\beta'})=C$ has at most $n/b$ vertices, this produces a tree decomposition of $\mathcal{Q}_{\alpha'\beta'}:=(T_{\alpha'\beta'},(C_x\mid x\in V(T_{\alpha'\beta'}))$ of height at most $\log_b(n/b)=\log_b n-1$.  By \ref{adhesion_pairs}, the root bag of $\mathcal{Q}_{\alpha'\beta'}$ has a tree-decomposition in which some bag contains the parent adhesions $A_{\alpha'}=B_{\alpha'}\cap B_x$ and $A_{\beta'}:=B_{\beta'}\cap B_y$.  We make the root bag of $T_{\alpha'\beta'}$ a child of $q_0$.  Doing this for each component $C$ of $T-S$ yields the desired tree-decomposition $\mathcal{Q}$ of $G$.
  
  Finally, to ensure that the resulting decomposition satisfies \ref{adhesion_pairs} we add $A_\alpha$ to $D_x$ for each $x\in V(T(q_0))$ on the path from $\alpha'$ to $\beta'$.  In particular $D_\beta$ contains $A_\alpha\cup A_\beta$.
\end{proof}


\subsection{Lower-Torso Specific}


Let $\mathcal{T}:=(T,(B_x\mid x\in V(T)))$ be a rooted tree-decomposition of a graph $G$. For each $x\in V(T)$ with parent adhesion $A_x$, define the \defin{lower-torso} $\mathdefin{\torso{G^-}{\mathcal{T},B_x}}:=\torso{G}{\mathcal{T},B_x}-A_x$.  For a child $y$ of a node $x$ of $T$, $A_y\setminus A_x$ is the \defin{lower-adhesion} between $x$ and $y$. The \defin{lower adhesion-width} of $\mathcal{T}$ is the size of the largest lower adhesion in $\mathcal{T}$.  The \defin{up-degree} of $G$ with respect to $\mathcal{T}$ is the the maximum over $x\in V(T)$ an $v\in B_x\setminus A_x$ of $|N_G(v)\cap A_x|$. 


% \begin{lem}
%   Let $G$ be a graph and let $\mathcal{T}$ be a tree-decomposition of $G$.  Then, for any $b>1$, there exists a tree-decomposition $\mathcal{Q}:=(Q,C_x\mid x\in V(T))$ with the following properties:
%   \begin{enumerate}[nosep,nolistsep,label=(\alph*)]
%     \item $\height(Q)\le\log_b n$.
%     \item 
%   \end{enumerate}
% \end{lem}



\begin{lem}
  Let $\mathcal{G}$ be a graph class that admits a weighted $(g_1,g_2,g_3)$ mixed labelling scheme.  Let $\mathcal{G}'$ be the graph class that consists of, for each $n\in\N$, every $n$-vertex graph $G$ that has a rooted tree-decomposition of size at most $b(n)$, up-degree at most $k(n)$, and whose lower-torsos belong to $\mathcal{G}$. Then $\mathcal{G}'$ admits a weighted $(g_1',g_2',g_3')$ mixed labelling scheme where
  \begin{align*}
    g_1'(n) & = g_1(n) + \Oh(k(n)\cdot(\log b(n)+g_2(n)+g_3(n)+\log\log n)) \\
    g_2'(n) & = g_2(n) + \Oh(\log b(n)+\log k(n)+\log\log n) \\
    g_3'(n) & = \Oh(k(n)\cdot (g_3(n)+\log k(n))+\log\log n) \enspace .
  \end{align*}
\end{lem}

\begin{proof}
  \pat{This proof fails to establish injectivity of $\mu$ on cliques.  I think the assumptions of the lemma are too weak to establish injectivity without increasing the clique label length to $\log b(n)$, which would make the lemma useless.  Back to working on the separator lemma above....}
  Let $A$ and $I$ be an adjacency tester and identity tester for a weighted $(g_1,g_2,g_3)$ mixed labelling for the class $\mathcal{G}$.    We will eventually define an adjacency tester $A'$ and identity tester $I'$ for $\mathcal{G}'$, but this is not practical without first describing how we obtain a labelling $(\mu,\kappa)$ of $(G^+,G,\omega)$ for $(A',I')$.
  
  Let $G^+$ be an $n$-vertex graph in $\mathcal{G}'$, let $G$ be a spanning subgraph of $G^+$, and let $\omega:V(G)\to\R^+$ be a weight function. By \cref{lem:nice_weights}, we may assume that $\omega$ is nice so that $\log\omega(G)\le \log n + O(1)$. Let $\mathcal{T}:=(T,(B_x\mid x\in V(T)))$ be a rooted tree-decomposition of $G^+$ witnessing that $G^+\in\mathcal{G}'$. For each $x\in V(T)$, let $A_x$ be the parent adhesion of $x$ in $\mathcal{T}$. For each node $y$ of $T$ with parent $x$, $|A_y\setminus A_x|\le k$, since $\mathcal{T}$ witnesses that $G\in\mathcal{G}'$. Let $k:=k(n)$ and let $b:=b(n)$.


  For each $x\in V(T)$, let $\omega(x):=\sum_{z\in V(T_x)}\omega(B_z\setminus A_z)$.  For each $x\in V(T)$, each vertex $u\in B_x\setminus A_x$ in the lower-torso of $x$, let $C_u$ be the set of children $y$ of $x$ such that $u$ is contained in $B_y\cap B_x\setminus A_x$.  Define
  \[  
    \delta(u) := \omega(u) + \sum_{y\in C_u}\omega(y) 
  \]
  Note that $\omega(y)$ appears in the sum for $\delta(u)$ exactly when $u$ is in the lower adhesion $A_y\setminus A_x$.  Since each lower-adhesion has size at most $k$, summing $\delta(u)$ over all $u\in B_x\setminus A_x$ implies that $\delta(B_x\setminus A_x)\le k\cdot\omega(x)$.

  Let $x_0,\ldots,x_r$ be a path from the root of $T$ to some node $x_r$.   For each $i\in\{0,\ldots,r-1\}$, let $K_{i}:=A_{x_{i+1}}\setminus A_{x_i}$. For each $i\in\{0,\ldots,r-1\}$ and each $u\in K_i$,  $\delta(u)\ge\omega(x_{i+1})$.  Therefore, 

  \begin{align*}
    &\quad \sum_{i=0}^{r-1}(\log \delta(B_{x_i}\setminus A_{x_i}) -\log\min_{u\in K_i}\delta(u)) 
    \le \sum_{i=0}^{r-1}(\log \delta(B_{x_i}\setminus A_{x_i}) -\log\omega(x_{i+1})) \\
    & \le \sum_{i=0}^{r-1}(\log (k\cdot\omega(x_i)) -\log\omega(x_{i+1})) 
    \le \log\omega(x_0)-\log\omega(x_r) + r\log k \\
    & = \log\omega(G)-\log\omega(x_r) + r\log k
  \end{align*}

  Furthermore, since $\omega(x_{i+1})\le \omega(x_i)$ for each $i$, for every $J\subseteq\{0,\ldots,r-1\}$,
  \begin{equation}
    \sum_{i\in J}(\log\delta(B_{x_i}\setminus A_{x_i})-\log\min_{u\in K_i}\delta(u))
    \le \log\omega(G)-\log\omega(x_r) + |J|\log k \label{telescoping_subset}
  \end{equation}

  For each $x\in V(T)$, let $(\mu_x,\kappa_x)$ be a weighted $(g_1,g_2,g_3)$ mixed labelling of $(\torso{G^+}{\mathcal{T},B_x}-A_x,G[B_x]-A_x,\delta)$ for $(A,I)$.  Let $\id:V(T)\to \{0,1\}^{\lceil\log b(n)\rceil}$ be an injective function that assigns a unique identifier to each node of $T$. For each $v\in V(G)$, let $h(v)$ be the home of $v$ in $\mathcal{T}$. That is, $h(v)$ is the unique $x\in V(T)$ such that $v\in B_x\setminus A_x$.

  \paragraph{\boldmath The vertex label $\mu(v)$:} For each $x\in V(T)$, and each $v\in B_x\setminus A_x$, we define
  \[
    \mu(v) := \langle \id(x),\mu_x(v),\tau(v),\alpha(v) \rangle \enspace .
  \]
  The parts $\id(x)$ and $\mu_x(v)$ are defined above. What remains is to describe $\tau(v)$ and $\alpha(v)$.  Let $x_0,\ldots,x_r$ be the path from the root $x_0$ of $T$ to the home $x_r=x$ of $v$. Let $u$ be a vertex in $N_G(v)\cap A_x$.  Then $u$ has a home $h(u)=x_i$ for some $i\in\{0,\ldots,r-1\}$. Since $u\in B_{x_i}\setminus A_{x_i}$ and $u\in A_{x_r}=A_x$, the vertex $u$ is in the lower-adhesion $K_{u}:=(B_{x_i}\cap B_{x_{i+1}})\setminus A_{x_i}$. In particular, $K_u$ is a non-empty clique in $\torso{G^+}{\mathcal{T},B_{h(u)}}-A_{h(u)}$ so $\mu_{h(u)}(K_u)$ and $\kappa_{h(u)}(K_u,u)$ are defined.  The bitstring $\tau(v)$ is used to store $K_u$ for each $u\in N_G(v)\cap A_x$:
  \[
    \tau(v) := \left\langle \langle\id(x_i), \mu_{x_i}((B_{x_i}\cap B_{x_{i+1}})\setminus A_{x_i}) \rangle \mid i\in\{0,\ldots,r-1\}, N_G(v)\cap B_{x_i}\setminus A_{x_i}\neq\emptyset \right\rangle \enspace .
  \]
  Thus, $\tau(v)$ is a sequence of at most $k$ pairs.
  For each $u\in N_G(v)\cap A_x$, exactly one of the pairs in $\tau(v)$ is $\langle\id(h(u)),\mu_{h(u)}(K_u)\rangle$.  For each $u\in N_G(v)\cap A_x$, let $i(u)$ denote the index in the sequence $\tau(v)$ of the pair $\langle\id(h(u)),\mu_{h(u)}(K_u)\rangle$.  Then
  \[
     \alpha(v) :=\left\langle \langle \bin(i(u)),\kappa_{h(u)}(K_u,u)\rangle \mid u\in N_G(v)\cap A_x \right\rangle  \enspace .
  \]
  
  \paragraph{\boldmath The vertex label length $|\mu(v)|$:}
  We now consider the length of the vertex label $\mu(v)$ for a vertex $v\in B_x\setminus A_x$. This label contains $\id(x)$ and $\id(h(u))$ for each $u\in N_G(v)\cap A_x$.  These identifiers have total length at most $(k+1)\lceil\log b\rceil$, so these contribute at most $\Oh(k\log b)$ bits to the length of $\mu(v)$.  In addition to this, $\alpha(v)$ contains up to $k$ pairs, each of which has length $\Oh(\log k + g_2(n))$.  The only remaining parts of $\mu(v)$ are $\mu_x(v)$ and $\{\mu_{h(u)}(K_u) \mid u\in N_G(v)\cap A_x\}$. By \eqref{telescoping_subset}, the lengths of these parts form a telescoping sum that is bounded by
  \begin{align*}
     &\quad \log\omega(G)-\log\omega(x)+\log\delta(B_x\setminus A_x)-\log\omega(v)+\Oh(k\cdot (g_3(n)+\log\log n)) \\
     & \le \log\omega(G)-\log\omega(v)+g_1(n)+\Oh(k\cdot (g_3(n)+\log\log n)).
  \end{align*}
  Thus $\mu(v)$ has total length at most $\log\omega(G)-\log\omega(v)+g_1(n)+\Oh(k\cdot(\log b+g_2(n)+g_3(n)+\log\log n))$, which satisfies the requirement for $g'_1(n)$.

  \paragraph{Adjacency testing:}  We now describe the adjacency tester $A'$ that takes two labels $\mu(v)$ and $\mu(w)$.  Suppose that $v\in B_x\setminus A_x$ and $w\in B_y\setminus A_y$ for some pair of node $x,y\in V(T)$.  Thus $\mu(v)=\langle \id(x),\mu_x(v),\tau(v),\alpha(v)\rangle$ and $\mu(w)=\langle \id(y),\mu_y(w),\tau(w),\alpha(w)\rangle$.  The adjacency tester first checks if $\id(x)=\id(y)$.  If $\id(x)=\id(y)$ then $x=y$, since $\id$ is injective. In this case $vw\in E(G)$ if and only if $vw\in E(G[B_x\setminus A_x])$, so $A'(\mu(v),\mu(w)):=A(\mu_x(v),\mu_y(w))$.

  If $\id(x)\neq \id(y)$ then the only way in which $v$ and $w$ can be adjacent in $G$ is if $w\in N_G(v)\cap A_x$ or $v\in N_G(w)\cap A_y$.  We now explain how to test if $v\in N_G(w)\cap A_y$. For each $u\in N_G(w)\cap A_y$, $\alpha(w)$ contains an entry $\langle i(u),\kappa_{h(u)}(K_u,u)\rangle$ where $K_u$ is a clique in $\torso{G^+}{\mathcal{T},B_{h(u)}}-A_{h(u)}$ that contains $u$ and $i(u)$ is the index of $\langle\id(h(u)),\mu_{h(u)}(K_u)\rangle$ in $\tau(w)$.  Thus, the adjacency tester $A'$ examines each entry $\langle i(u),\kappa_{h(u)}(K_u,u)\rangle$.  For each entry, it checks if $\id(h(u))=\id(x)$.  If not, then $h(u)\neq x$, $u\in B_{h(u)}\setminus A_{h(u)}$ and $v\in B_x\setminus A_x$, so $u\neq v$.  Otherwise, $x=h(u)$ and the tester checks if $I(\mu_{h(u)}(K_u),\kappa_{h(u)}(K_u,u),\mu_x(v))=1$.  If so, then $u=v$ so $v\in N_G(w)\cap A_y$.  In particular, $v\in N_G(w)$ so $A'(\mu(v),\mu(w)):=1$.  By doing this for each entry $\langle i(u),\kappa_{h(u)}(K_u,u)\rangle$ in $\alpha(w)$ the adjacency tester either returns $1$ because $v\in N_G(w)\cap A_y$ or determines that $v\not\in N_G(w)\cap A_y$.  The tester then repeats this process with the roles of $v$ and $w$ reversed and either returns $1$ because $w\in N_G(v)\cap A_x$ or determines that $w\not\in N_G(v)\cap A_x$.  In the latter case, $A'(\mu(v),\mu(w)):=0$ since $vw\not\in E(G[B_x\setminus A_x])$, $v\not\in N_G(w)\cap A_y$ and $w\not\in N_G(v)\cap A_x$.

  \paragraph{The clique label $\mu(K)$:}
  Let $G^\star:=\bigcup_{x\in V(T)}\torso{G^+}{\mathcal{T},B_x}$.  Note that $G^\star$ is a supergraph of $G^+$ since $\torso{G^+}{\mathcal{T},B_x}$ is a supergraph of $G^+[B_x]$ for each $x\in V(T)$.

  Let $K$ be a clique in $G^+$.  Consider the set $h(K):=\{x\in V(T):K\cap B_x\setminus A_x\neq\emptyset\}$.  Since $K$ is a clique in $G^+$, there is a single path $x_0,\ldots,x_r$ from the root $x_0$ of $T$ to the deepest node $x_r=z$ such that $K\cap B_{z}\setminus A_z$ is non-empty.  Since $K$ is a clique, a vertex in $K\cap B_{z}\setminus A_z$ has at least $|h(K)|-1$ neighbours in $A_z$.  Since $\mathcal{T}$ has up-degree at most $k$, this implies that $|h(K)|\le k+1$.
  
  For each $x\in h(K)$, $K\cap B_x\setminus A_x$ is a non-empty clique in $\torso{G^+}{\mathcal{T},B_x}-A_x$, so $\mu_x(K\cap B_x\setminus A_x)$ is defined.  Define the clique label
  \[
    \mu(K)
      :=\left\langle \langle\mu_x(K\cap B_x\setminus A_x)\rangle
          \mid x\in h(K)
        \right\rangle
  \]
  By \eqref{telescoping_subset}, the lengths of the at most $k+1$ clique labels in $\mu(K)$ form a telescoping series, so
  \begin{align*}
    \sum_{x\in h(K)}|\mu_{x}(K\cap B_x\setminus A_x)|
    & \le 
    \log\omega(G)-\log\min_{u\in K\cap B_z\setminus A_z}(\omega(u)) + \Oh(k\cdot(g_3(n)+\log k)) \\
    & \le \log\omega(G)-\log\min_{u\in K}(\omega(u))
    +\Oh(k\cdot(g_3(n)+\log k))
  \end{align*}
  Thus, $|\mu(K)|\le \log\omega(G)-\log\min_{u\in K}(\omega(u))
    +\Oh(k\cdot(g_3(n)+\log k+\log\log n))
    $, which satisfies the requirement for $g'_3(n)$.

  \paragraph{\boldmath The local identifier $\kappa(K,u)$:} For each $u\in K$, let $i(u)$ be the index of  $\mu_{h(u)}(K\cap B_{h(u)}\setminus A_{h(u)})$ in the sequence $\mu(K)$.  For each $u\in K$, define the local identifier
  \[ 
    \kappa(K,u):=\langle \bin(i(u)),\id(h(u)),\kappa_{h(u)}(K\cap B_{h(u)}\setminus A_{h(u)}, u) \rangle
  \]
  The length of $\kappa(K,u)$ is $g_2(n)+\Oh(\log b+\log k+\log\log n)$, which satisfies the requirement for $g'_2(n)$.

  \paragraph{Identity Testing:}  We now describe the operation of the identity tester $I'$ which takes as input the label $\mu(K)$ of a clique $K$ in $G^+$, the local identifier $\kappa(K,u)=\langle\bin(i(u)),\id(h(u)),\kappa_{h(u)}(K\cap B_{h(u)}\setminus A_{h(u)}, u)\rangle$ of a vertex $u\in K$, and the label $\mu(v)=\langle \id(x),\mu_x(v),\alpha(v)\rangle$ of a vertex $v\in B_x\setminus A_x$ for some $x\in V(T)$.  The identity tester $I'$ first checks if $\id(h(u))=\id(x)$.  If not, then $h(u)\neq x$, so $u\in B_{h(u)}\setminus A_{h(u)}$ and $v\in B_x\setminus A_x$, so $v\neq u$ and $I'(\mu(K),\kappa(K,u),\mu(v)):=0$. If $\id(h(u))=\id(x)$ then $h(u)=x$, since $\id$ is injective.  In this case, the identity tester uses $i(u)$ to extract $\mu_{h(u)}(K\cap B_{h(u)}\setminus A_{h(u)})$ from $\mu(K)$.  Then the identity tester returns $I'(\mu(K),\kappa(K,u),\mu(v)):=I(\mu_{h(u)}(K\cap B_{h(u)}\setminus A_{h(u)}), \kappa_{h(u)}(K\cap B_{h(u)}\setminus A_{h(u)}, u), \mu_{x}(v))$.
\end{proof}






% SAVE E4cRTXgmT7EGUyz!!!
S
